% Part: sets-functions-relations
% Chapter: functions
% Section: isomorphisms

\documentclass[../../../include/open-logic-section]{subfiles}


\begin{document}

\olfileid{sfr}{fun}{iso}
\olsection{التشاكل}

\begin{explain}
\emph{التشاكل} هو تقابل يحفظ بنية المجموعتين اللتين يربط بينهما، حيث تتمثل
البنية في العلاقات القائمة بين !!{element}s المجموعتين. انظر إلى
المجموعتين $X=\{1,2,3\}$ و$Y=\{4,5,6\}$. تنظم كلًّا منهما علاقات
التالي، وأصغر من، وأكبر من. ويكون التشاكل بينهما !!a{bijection} يحفظ
هذه البنى. ومن ثم، إذا كانت الدالة $f \colon X \to Y$ !!a{bijective}،
فإنها تكون تشاكلًا متى تحقق، لكل $i,j\in X$، أن $i<j$ إذا وفقط إذا
كان $f(i)<f(j)$، وأن $i>j$ إذا وفقط إذا كان $f(i)>f(j)$، وأن $j$ هو
تالي $i$ إذا وفقط إذا كان $f(j)$ تالي $f(i)$.
\end{explain}

\begin{defn}[التشاكل]
لتكن $U$ الزوج $\langle X, R\rangle$، ولتكن $V$ الزوج $\langle
Y, S\rangle$، بحيث تكون $X$ و$Y$ مجموعتين وتكون $R$ و$S$ علاقتين
على $X$ و$Y$، على الترتيب. وليكن $f$ !!{bijection} من $X$ إلى
$Y$. يكون $f$ \emph{تشاكلًا} من $U$ إلى $V$ إذا وفقط إذا كان يحفظ
البنية العلائقية؛ أي إنه، لكل $x_{1}$ و$x_{2}$ في $X$،
$\tuple{x_1,x_2} \in R$ إذا وفقط إذا كان $\tuple{f(x_1),f(x_2)} \in S$.
\end{defn}

\begin{ex}
انظر إلى المجموعتين $X=\{1,2,3\}$ و$Y=\{4,5,6\}$، وإلى علاقتي
أصغر من وأكبر من. إن الدالة $f\colon X \to Y$ المعرّفة بـ
$f(x) = 7-x$ تشاكل بين $\tuple{X,<}$ و$\tuple{Y,>}$.
\end{ex}

\end{document}
